\subsection{Chegada mais cedo e variação da espera}\label{sec:F-ESPERA-2}
Em $G_T$, o Dijkstra rotula paradas pelo instante de chegada. A busca binária
encontra a primeira partida viável após descartar viagens ultrapassadas.
Com $k$ horários por segmento, a consulta custa
$O((V+E)\log(E+2)+E\log(k+1))$. O Monte Carlo repete a consulta para saídas
uniformes e registra quais alcançam o destino.

\begin{center}\small
\begin{tabular}{@{}ll@{}}
 & \textbf{Espera no trajeto (min)} \\
Quadro de horários & \rule{0.374cm}{0.7em}\enspace 1,87 \\
Estimativa $H/2$ & \rule{3.068cm}{0.7em}\enspace 15,34 \\
\end{tabular}
\end{center}
No GTFS de BH, entre as paradas 14788851 e 14792190 às 08h do serviço de
04/11, o total foi de 23,67 min pelo quadro e 40,24 min por $H/2$. Cada
modelo escolhe seu trajeto. O quadro usa partidas programadas e admite
conexões no mesmo segundo.
